\documentclass[a4paper]{article}

% Basic packages
\usepackage[fontset=fandol]{ctex}
\usepackage{amsmath, amssymb}
\usepackage{graphicx}
\usepackage[margin=2.5cm]{geometry}
\usepackage[most]{tcolorbox}
\usepackage{etoolbox}
\usepackage{listings}
\usepackage{hyperref}
\usepackage{booktabs}
\usepackage{subcaption}
\usepackage{float}
\usepackage{tikz}
\IfFileExists{pgfplots.sty}{
    \usepackage{pgfplots}
    \pgfplotsset{compat=1.18}
}{}

% Highlight boxes
\newtcolorbox{knowledgebox}[1]{
    enhanced, colback=blue!5!white, colframe=blue!75!black, colbacktitle=blue!75!black,
    coltitle=white, fonttitle=\bfseries, title=#1, attach boxed title to top left={yshift=-2mm, xshift=2mm},
    boxrule=1pt, sharp corners
}

\newtcolorbox{importantbox}[1]{
    enhanced, colback=yellow!10!white, colframe=yellow!80!black, colbacktitle=yellow!80!black,
    coltitle=black, fonttitle=\bfseries, title=#1, sharp corners
}

\newtcolorbox{warningbox}[1]{
    enhanced, colback=red!5!white, colframe=red!75!black, colbacktitle=red!75!black,
    coltitle=white, fonttitle=\bfseries, title=#1, sharp corners
}

% Listings
\lstset{
    language=Python,
    basicstyle=\ttfamily\small,
    keywordstyle=\color{blue},
    stringstyle=\color{red!60!black},
    commentstyle=\color{green!60!black},
    breaklines=true,
    frame=single,
    numbers=left,
    numberstyle=\tiny\color{gray},
    captionpos=b,
    extendedchars=false
}

% Front-page metadata
\newcommand{\notetitle}{CS231N 2025 第3讲：正则化与优化}
\newcommand{\noteauthors}{Codex Harness Worker 1}
\newcommand{\notedate}{2026-04-15}
\newcommand{\videochannel}{Stanford Online}
\newcommand{\videopublishdate}{2025-09-02}
\newcommand{\videoduration}{1:08:39}
\newcommand{\videourl}{https://www.youtube.com/watch?v=dyNGd06MWn4}
\newcommand{\videocoverpath}{cover.jpg}

\begin{document}

\begin{titlepage}
\centering
{\Large 课程笔记\par}
\vspace{1.2cm}
{\huge\bfseries \notetitle\par}
\vspace{0.8cm}
{\large \noteauthors\par}
\vspace{0.3cm}
{\large \notedate\par}
\vspace{1.2cm}

\ifdefempty{\videocoverpath}{
{\small 请在生成笔记时填入视频封面图的本地路径。\par}
}{
\includegraphics[width=0.82\textwidth,height=0.45\textheight,keepaspectratio]{\videocoverpath}\par
}

\vfill
\begin{tcolorbox}[width=0.9\textwidth, colback=black!2!white, colframe=black!60, sharp corners]
\textbf{视频作者/频道}：\videochannel\par
\textbf{发布日期}：\videopublishdate\par
\textbf{视频时长}：\videoduration\par
\textbf{视频链接}：
\ifdefempty{\videourl}{
未填写
}{
\href{\videourl}{\nolinkurl{\videourl}}
}
\end{tcolorbox}
\end{titlepage}

\tableofcontents
\newpage

%% --- 正文内容开始 --- %%


\section{目标函数：经验风险最小化加上复杂度控制}
上一讲已经把图像分类写成“分数函数 + 损失函数”的形式，这一讲则进一步把训练目标写成完整优化问题。课程最强调的统一表达是：
\[
L(W) = \frac{1}{N} \sum_{i=1}^{N} L_i\bigl(f(x_i; W), y_i\bigr) + \lambda R(W).
\]
其中，$L_i$ 是第 $i$ 个样本的数据损失，$R(W)$ 是正则项（regularizer），$\lambda$ 控制正则强度。这个式子看似简单，却是后续几乎所有深度学习训练问题的原型。

符号上，$N$ 表示训练样本数，$f(x_i;W)$ 是由参数 $W$ 给出的预测分数，$y_i$ 是真实标签，$R(W)$ 则衡量模型复杂度。教材化理解时，务必把它拆成两个问题：第一，模型是否足够贴合训练数据；第二，模型是否因为过度复杂而失去泛化能力。正则化正是在第二个问题上发挥作用。

\begin{figure}[H]
\centering
\includegraphics[width=0.92\textwidth]{pdf_pages/page-20.png}
\caption{损失由数据项与正则项组成，正则化强度控制“拟合训练集”与“保持参数简单”之间的平衡。}
\end{figure}

正则化与数据损失共同定义了真正要最小化的目标函数。
深度学习训练从来不是单目标，而是拟合与泛化之间的权衡。
\subsection{本章小结}

\section{为什么要正则化：泛化而非死记硬背}
正则化（regularization）最常见的直觉，是“防止模型把训练集背下来”。但更精确地说，正则化是在向模型注入偏好：偏好更小的权重、更平滑的函数、更少的自由度，或者更稳定的预测行为。对于线性分类器，这常常体现为 $L_2$ 范数惩罚；对于神经网络，则还会体现在 dropout、数据增强、早停（early stopping）等更丰富的策略上。

课程在这里特别强调一个容易被忽略的事实：较低的训练误差并不保证较低的测试误差。真正的目标是测试分布上的表现，而不是训练分布上的记忆能力。因此，判断一个优化策略是否有效，必须看它在验证集上的行为，而不是只盯着训练 loss 曲线。

正则化不是训练失败时才加的补丁，而是现代模型设计的组成部分。
泛化能力必须通过验证集与测试集来检验，而不能从训练误差直接推断。
\subsection{本章小结}

\section{如何得到梯度：数值法与解析法}
一旦目标函数写出来，下一步就是求梯度。课程先用有限差分（finite differences）解释数值梯度：
\[
\frac{\partial L}{\partial w_k} \approx \frac{L(w_k + h) - L(w_k - h)}{2h}.
\]
其中，$h$ 是一个很小的扰动，$w_k$ 是某个参数。这个公式的意义在于：即便你不会手推导数，也能通过“微小扰动前后的 loss 变化”近似得到梯度。

但数值梯度只适合做 gradient check。它的优点是通用、容易实现；缺点是极慢，而且会受数值误差影响。真正训练模型时，我们必须依赖解析梯度（analytic gradient），也就是用微积分直接推导每个参数对 loss 的偏导数。后面反向传播的全部价值，就体现在“如何高效得到解析梯度”。

\begin{figure}[H]
\centering
\includegraphics[width=0.92\textwidth]{pdf_pages/page-40.png}
\caption{数值梯度示意：通过有限差分近似偏导数。它常用于 gradient check，但训练时过慢。}
\end{figure}

有限差分是检查梯度是否写对的工具，不是训练大模型的方法。
训练时必须依赖可组合的解析梯度，这直接引向反向传播。
\subsection{本章小结}

\section{从全量梯度下降到小批量 SGD}
理论上，最直接的更新规则是全量梯度下降：
\[
W \leftarrow W - \eta \nabla_W L(W),
\]
其中，$\eta$ 是学习率（learning rate），$\nabla_W L(W)$ 是对全体训练集求出的梯度。但在大规模数据下，全量梯度每一步都要扫完整个数据集，代价太高，因此实际训练通常使用小批量随机梯度下降（mini-batch SGD）。

mini-batch SGD 用一个小批样本近似整体梯度。这样得到的更新方向带有噪声，但这种噪声并不总是坏事：它能降低每步成本，也常常帮助模型逃离某些坏的局部结构。课程在这里想传达的关键信息是：优化不是只看“每一步有多准”，还要看“在固定算力预算下能走多少步”。

SGD 把每步更新从“极准但极贵”改成“略噪但高效”，这是深度学习可扩展性的关键。
学习率决定每步走多远，而批大小决定梯度估计有多噪。
\subsection{本章小结}

\section{从 SGD 到 Momentum，再到 Adam}
纯 SGD 的一个问题是，在狭长谷地中常常左右震荡、前进缓慢。动量法（momentum）通过累积历史梯度来缓解这个问题：
\[
 v_t = \mu v_{t-1} - \eta \nabla_W L(W_t), \qquad W_{t+1} = W_t + v_t.
\]
其中，$v_t$ 是速度项，$\mu$ 是动量系数。它的直觉很像在优化地形上滚动一个有惯性的球：沿一致方向会不断加速，来回震荡方向则会相互抵消。

\begin{figure}[H]
\centering
\includegraphics[width=0.92\textwidth]{pdf_pages/page-60.png}
\caption{动量法把历史梯度累积成“速度”，从而在狭长谷地中减少震荡、加速沿主方向前进。}
\end{figure}

自适应方法则进一步试图为不同参数分配不同步长。以 Adam 为例，可写成
\[
 m_t = \beta_1 m_{t-1} + (1-\beta_1) g_t, \qquad
 v_t = \beta_2 v_{t-1} + (1-\beta_2) g_t^2,
\]
\[
 W_{t+1} = W_t - \eta \frac{\hat{m}_t}{\sqrt{\hat{v}_t} + \epsilon}.
\]
其中，$g_t$ 是当前梯度，$m_t$ 与 $v_t$ 分别是一阶矩和二阶矩的指数滑动平均，$\hat{m}_t, \hat{v}_t$ 是偏差校正后的估计，$\epsilon$ 用于数值稳定。这个公式说明 Adam 通过参数级步长调整，适应不同尺度和稀疏度的更新。

\begin{figure}[H]
\centering
\includegraphics[width=0.92\textwidth]{pdf_pages/page-70.png}
\caption{Adam 结合一阶矩与二阶矩的指数滑动平均，在不同参数尺度上自适应调整步长。}
\end{figure}

Momentum 通过历史速度减少震荡、提升主方向推进效率。
Adam 则通过一阶矩和二阶矩估计，实现参数级自适应步长。
\subsection{本章小结}

\section{学习率日程与二阶方法的角色}
即使优化器固定，学习率日程（learning rate schedule）仍然决定了训练能否真正收敛。课程展示了 step decay 等思路，其共同逻辑是：训练早期需要较大学习率快速探索，训练后期则需要更小步长把参数稳定压到更优区域。

\begin{figure}[H]
\centering
\includegraphics[width=0.92\textwidth]{pdf_pages/page-80.png}
\caption{学习率衰减策略决定了训练后期能否稳定收敛；它通常与动量和批大小共同作用。}
\end{figure}

在理论上，二阶方法利用 Hessian 信息能够更精确描述局部曲率，但在深度网络中往往因存储和计算开销过高而难以直接使用。因此实践中的主流方案仍是“一阶优化器 + 合理学习率调度 + 正则化”。

学习率是训练动力学中最敏感的超参数之一，日程设计常常和优化器本身同等重要。
深度学习实践仍以一阶方法为主，二阶方法更多作为理解曲率和局部几何的参考。
\subsection{本章小结}

\section{总结与延伸}
Lecture 3 解决的是一个朴素但决定性的问题：写下目标函数之后，怎样把参数真正学出来，而且学出来的模型还能泛化。为此，课程把训练拆成四个层次：目标定义、正则化、梯度求法、更新策略。只有四者协同，训练过程才是完整的。

从教材角度看，这一讲的地位极高，因为后续神经网络、卷积网络、自监督学习乃至大模型训练，都没有跳出这个框架。目标函数可能换，模型结构可能变，但“正则化 + 梯度 + 优化器 + 学习率日程”的组合始终存在。下一讲进入神经网络与反向传播时，读者应该把本讲内容当作它的优化背景板。

\subsection{拓展阅读}
\begin{itemize}
\item \href{https://cs231n.github.io/optimization-1/}{Optimization}
\end{itemize}


%% --- 正文内容结束 --- %%

\end{document}
